\documentclass[11pt]{article}
\usepackage[T1]{fontenc}
\usepackage{lmodern}
\usepackage[margin=1in]{geometry}
\usepackage{amsmath,amssymb,amsthm,mathtools,microtype,enumitem,setspace,xcolor}
\usepackage[colorlinks=true,linkcolor=blue!45!black,citecolor=blue!45!black,
            urlcolor=blue!45!black,breaklinks=true]{hyperref}
\onehalfspacing
\numberwithin{equation}{section}
\newtheorem{theorem}{Theorem}[section]
\newtheorem{proposition}[theorem]{Proposition}
\newtheorem{lemma}[theorem]{Lemma}
\newtheorem{corollary}[theorem]{Corollary}
\theoremstyle{remark}
\newtheorem{remark}[theorem]{Remark}
\newcommand{\E}{\mathbb E}
\newcommand{\T}{\mathcal T}
\newcommand{\Unif}{\operatorname{Unif}}
\newcommand{\depth}{\operatorname{depth}}
\newcommand{\Wone}{W_1}
\setlist[enumerate]{itemsep=2pt,topsep=4pt}
\title{Optimal fully dynamic matching on the line}
\ifdefined\AnonymousSubmission
  \author{}
  \hypersetup{pdftitle={Optimal fully dynamic matching on the line},pdfauthor={}}
\else
  \author{Yash Kanoria\thanks{Decision, Risk, and Operations Division,
  Columbia Business School, Columbia University, New York, NY 10027, USA.
  Email: \href{mailto:ykanoria@columbia.edu}{ykanoria@columbia.edu}.}}
  \hypersetup{pdftitle={Optimal fully dynamic matching on the line},pdfauthor={Yash Kanoria}}
\fi
\date{}
\begin{document}
\maketitle
\begin{abstract}
We study online matching with a fixed inventory of \(m\) supply points on
the unit interval. Each uniform demand is matched immediately to an
available supply, which departs and is replaced by an independent uniform
point. We construct a policy with long-run expected matching distance
\(O(\log m/m)\), matching the known lower bound. The same order holds for
the long-run root-mean-square distance and, after initialization, for average cost
over every horizon of at least \(2m^2\) periods. The policy assigns deletion
probabilities recursively on a dyadic tree and implements them by quantile
transport. Its analysis combines a harmonic inventory potential with a
quadratic Bellman inequality. Tree symmetry cancels interactions between
spatial scales, while the Bellman inequality controls both transport
discrepancy and the curvature loss in the potential drift.
\end{abstract}
\noindent\textbf{Keywords:} dynamic matching; optimal transport; stochastic
control; Lyapunov drift; multiscale analysis.

\section{Introduction}

Spatial matching with replenishment requires a balance between current
service cost and future availability. Matching each demand to a nearby
supply saves distance immediately, but the resulting inventory can become
depleted in regions that future demands will visit. In one dimension,
controlling this depletion simultaneously at many spatial scales is the
central difficulty.

We consider a model with \(m\) available supplies, independent uniform
demand arrivals, and independent uniform replenishment after every match.
Kanoria \cite{Kanoria} established a lower bound of order \(\log m/m\)
and an upper bound of order \((\log m)^2/m\) for this model. The lower
bound relates dynamic matching on the line to static random matching in
two dimensions, with time playing the role of the second coordinate, and
uses the theorem of Ajtai, Koml\'os and Tusn\'ady \cite{AKT}. We close
this logarithmic gap by constructing a policy with cost \(O(\log m/m)\).

The construction separates inventory control from spatial assignment.
Conditional on the inventory, it first specifies the probability of
deleting a supply from each cell of a dyadic partition. Monotone transport
then couples these probabilities to the arriving demand. The deletion
probabilities determine the inventory dynamics; their coupling to demand
determines the immediate distance. This separation allows the inventory
feedback to be designed for a tractable drift calculation.

\subsection{Model and main result}

Fix \(m\ge2\). The initial inventory \(S_0\) is a multiset of \(m\) points
in \([0,1]\). Let \((U_t)_{t\ge1}\) and \((A_t)_{t\ge1}\) be mutually
independent iid \(\Unif[0,1]\) sequences. In period \(t\), the operator
observes \(U_t\), chooses \(X_t\in S_{t-1}\), incurs cost
\(C_t=|U_t-X_t|\), and deletes \(X_t\). The replenishment \(A_t\) is
then observed and inserted. A policy may use the current demand and the
history, but not future demands or replenishments. Write
\[
J_T=\frac1T\sum_{t=1}^T\E C_t,
\qquad J_\infty=\limsup_{T\to\infty}J_T.
\]
All logarithms are natural unless a base is indicated.

\begin{theorem}\label{thm:main}
There is an online policy and a universal constant \(C<\infty\) such
that, for every \(m\ge2\) and every initial inventory,
\begin{enumerate}[label=(\roman*)]
\item the policy satisfies
\[
\limsup_{t\to\infty}(\E C_t^2)^{1/2}
\le C\frac{\log(m+1)}m;
\]
\item with an initial phase that replaces the \(m\) original supplies,
its average cost satisfies
\[
J_T\le C\frac{\log(m+1)}m
\qquad\text{for every }T\ge2m^2.
\]
\end{enumerate}
The policy can be implemented using \(O(\log m)\) operations per period
and \(O(m)\) memory in the real-arithmetic model.
\end{theorem}

The initial replacement phase consists of matching to one original supply
in each of the first \(m\) periods and retaining all replenishments.
It produces iid uniform inventory at total cost at most \(m\).
The constants in the proof are explicit but are not optimized.

\begin{corollary}\label{cor:sharp}
For every \(m\ge2\), every initial inventory, and every \(T\ge2m^2\),
\[
\inf_\alpha J_T(\alpha)=\Theta\!\left(\frac{\log m}{m}\right),
\]
with universal implied constants. The same order holds for the optimal
long-run expected cost.
\end{corollary}
\begin{proof}
Apply Theorem~\ref{thm:main} and the lower bound in
Kanoria \cite[Theorem~3]{Kanoria}.
\end{proof}

\subsection{Related work and proof outline}

Shor \cite{Shor} used hierarchical inventory control in stochastic packing
and matching problems. These constructions motivate multiscale feedback,
although the present policy and its drift estimate are different.
Hierarchical Greedy, analyzed by Kanoria \cite{Kanoria}, gives the previous
\(O((\log m)^2/m)\) upper bound in the fully dynamic one-dimensional
model. Our result concerns the optimal achievable cost; the sharp
performance of Hierarchical Greedy remains a separate question.

The separation of deletion probabilities from their spatial implementation
has a close precedent in Holden, Peres and Zhai \cite{HPZ}. Their
allocation-based online matching construction gives every available supply
equal deletion probability, preserving the iid distribution of the
survivors. In the fully dynamic model, preserving iid inventory is
insufficient for the desired bound \cite[Section~7.1]{Kanoria}.
Here the deletion probabilities depend on the inventory, and the analysis
controls the resulting dependent count distribution.
For the semidynamic model, in which supplies are present initially and
are not replenished, Balkanski, Faenza and P\'erivier \cite{BFP} establish
constant competitiveness of Greedy when supplies and demands are iid
uniform. Replenishment makes persistent spatial depletion a distinct
control problem.

The main analytic step is a local quadratic Bellman inequality, chosen
together with the policy. Summed down the tree, it controls squared
deletion imbalances by the restoring term of a harmonic inventory
potential. The curvature loss in the potential drift is bounded by one
common rate energy per level. A regularization parameter proportional to
the tree depth absorbs this loss. Finally, an integrated Haar expansion
relates the deletion imbalances to transport distance. Its cross terms
vanish in expectation by symmetry of the inventory dynamics; this
cancellation is essential for the logarithmic order.

Section~\ref{sec:policy} defines the policy and proves its elementary
properties. Section~\ref{sec:energy} develops the drift estimate, and
Section~\ref{sec:transport} converts it to matching-cost bounds.
Section~\ref{sec:heuristic} explains the design choices after the policy
and potential have been introduced. Section~\ref{sec:extensions} discusses
further directions and limitations. The local Bellman calculation is given in
Appendix~\ref{app:local-proof}.

\section{The matching policy}\label{sec:policy}

Set
\begin{equation}\label{eq:parameters}
a=2000\lceil\log_2(m+1)\rceil,\qquad
n=2^L=2^{\lfloor\log_2\max\{1,m/a\}\rfloor}.
\end{equation}
Thus
\begin{equation}\label{eq:parameter-bounds}
\frac1n\le\frac{2a}{m},\qquad a\ge2000L.
\end{equation}


Partition \([0,1]\) into \(n=2^L\) equal dyadic leaves and let \(\T\)
be the complete binary tree of their intervals, with root \(o\).
Let \(\T_\ell\) denote the nodes at depth \(\ell\). Use half-open cells,
except for the rightmost cell, so each supply belongs to exactly one leaf.
A node \(v\) has interval
\(I_v\), length \(p_v=2^{-\depth(v)}\), and inventory \(N_v\).

Starting with \(q_{\mathrm{root}}=1\), assign deletion masses recursively.
At a nonempty internal node abbreviate its length and mass by \(p,q\).
Temporarily order its child inventories as \(x\ge y\), put
\(N=x+y\) and \(\Delta=x-y\), and, when \(y>0\), set
\begin{equation}\label{eq:policy}
b=\min\left\{
\frac{2aq\Delta}{xy},
\frac{q\Delta}{N},
q-\frac{2py}{2y+a}
\right\}.
\end{equation}
Give deletion mass \((q+b)/2\) to the larger child and \((q-b)/2\) to
the smaller, and then restore their spatial labels. If \(y=0\), send all
mass to the occupied child. If \(x=y>0\), split equally. Empty nodes and
their descendants receive mass zero.

The three terms in \eqref{eq:policy} will be called the \emph{feedback},
\emph{uniform-deletion cap}, and \emph{rate-floor cap}. The feedback
favors the more populated child. The second term keeps that child's
average deletion rate per supply below its sibling's; the third enforces
a lower bound \(p/(2y+a)\) on the smaller child's rate. Lemma
\ref{lem:invariants} proves that they define valid probabilities and that
every occupied child receives positive mass.

\subsection{Quantile implementation}

Spread each leaf mass \(q_i\) uniformly over its leaf, producing a
probability measure \(Q\) with distribution function \(F_Q\). Given demand
\(U\), let \(Y=F_Q^{-1}(U)\), choose the leaf containing \(Y\), and match
to a fixed representative supply of that leaf. Representatives can be
maintained in linked lists. The choice of representative has no effect
on the count dynamics or on our cost bounds.
Equivalently, route \(U\) recursively using the conditional child masses,
rescaling the residual quantile at each split. Only the visited path is
evaluated. After deletion and replenishment, update the two affected
count paths. The cost is \(O(L+1)\) operations per period and \(O(m+n)\)
memory.

\begin{proposition}[Separation by quantile transport]\label{prop:quantile}
For any fixed state-dependent deletion probabilities on ordered supplies,
monotone quantile assignment preserves the inventory transition law and
minimizes immediate expected distance among assignments with those deletion
probabilities. For the leaf implementation above,
\begin{equation}\label{eq:leaf-cost}
\E[C\mid S]
\le\frac1n+\int_0^1|F_Q(z)-z|\,dz.
\end{equation}
\end{proposition}
\begin{proof}
The quantile intervals have exactly the prescribed masses, so the deleted
label and hence the next-inventory law are unchanged. Monotone coupling is
optimal for absolute distance on the line and has cost
\(\Wone(\Unif[0,1],Q)=\int_0^1|F_Q(z)-z|\,dz\). The actual supply selected
inside the leaf lies within one leaf width of the auxiliary point \(Y\),
which proves \eqref{eq:leaf-cost}.
\end{proof}

\subsection{Feasibility and rate invariants}\label{sec:structure}

For a nonempty node define its average deletion rate per supply by \(h=q/N\), and
at every node define
\begin{equation}\label{eq:htz}
t=\frac{p-q}{a},
\qquad
Z=\frac{p}{N+a/2}.
\end{equation}
The variables \(t\) and \(Z\) are defined by \eqref{eq:htz} even when
\(N=q=0\). We extend \(h\) analytically to empty nodes as follows. At an
empty child of a nonempty parent set
\begin{equation}\label{eq:empty-rate}
h_{\mathrm{empty}}
=\max\{h_{\mathrm{parent}},p_{\mathrm{parent}}/a\}.
\end{equation}
At an empty parent, pass its auxiliary rate to both children. These rates
are used only in the analysis; no deletion mass is assigned to an empty
node.

For a parent whose children are ordered as \(x\ge y\), write
\(h_+=(q+b)/(2x)\) and \(h_-=(q-b)/(2y)\) when both counts are positive.
Let \(b_v=q_{v_{\mathrm L}}-q_{v_{\mathrm R}}\) use the fixed spatial
left--right labels.

\begin{lemma}[Invariant domain and rate energy]\label{lem:invariants}
At every node,
\begin{equation}\label{eq:invariant-domain}
t\le\frac h2,
\qquad Z\le h.
\end{equation}
For a positive-count split, all three terms in \eqref{eq:policy} are
nonnegative, \(0\le b<q\), and every occupied child receives positive
mass. Moreover, if
\begin{equation}\label{eq:H-G}
H_\ell=\sum_{v\in\T_\ell}p_vh_v^2,
\qquad
G=\sum_{v\in\T\setminus\T_L}p_vb_v^2,
\end{equation}
then
\begin{equation}\label{eq:H-monotone}
H_0=m^{-2}\le H_1\le\cdots\le H_L.
\end{equation}
\end{lemma}
\begin{proof}
At the root, \(t=0\) and \(h=1/m\). For a nonempty node the two
inequalities in \eqref{eq:invariant-domain} are equivalent: both reduce to
\[
p\le q\left(1+\frac{a}{2N}\right).
\]
Suppose this holds at a nonempty parent. The feedback and uniform-deletion
cap are nonnegative, and the parent invariant gives
\[
q-\frac{2py}{2y+a}\ge\frac{ahx}{2y+a}>0.
\]
When \(y>0\), the uniform cap is strictly less than \(q\), so
\(0\le b<q\) and both child masses are positive.
The rate-floor cap gives the
smaller child mass
\[
q_-\ge\frac{py}{2y+a},
\]
which is exactly \(Z_-\le h_-\). For the larger child, \(q_+\ge q/2\)
and
\[
p\le q\left(1+\frac a{2N}\right)
\le q\left(1+\frac a{2x}\right),
\]
which gives \(Z_+\le h_+\). If the smaller child is empty,
\eqref{eq:empty-rate} gives \(h_-=p/a=Z_-\) or a larger value; below an
empty parent the conclusion is inherited.

The uniform-deletion cap is equivalent to \(h_+\le h_-\). The
parent rate is the inventory-weighted average
\(h=(xh_++yh_-)/(x+y)\), so
\[
\frac{h_++h_-}{2}-h
=\frac{x-y}{2(x+y)}(h_--h_+)\ge0.
\]
Jensen's inequality therefore gives
\[
\frac{h_+^2+h_-^2}{2}\ge h^2.
\]
The auxiliary rates give the same inequality at empty children. Multiplying
by \(p\), and observing that each child length is \(p/2\), proves
\eqref{eq:H-monotone} after summing over a level.
\end{proof}

\subsection{Symmetry of the count process}

\begin{lemma}[Count-chain symmetry]\label{lem:symmetry}
The finite chain of leaf counts is irreducible and aperiodic and hence has
a unique stationary law \(\pi\). This law is invariant under every
automorphism of the rooted binary tree. An invariant initial count law
remains invariant at every deterministic time.
\end{lemma}
\begin{proof}
Every occupied leaf has positive deletion probability, so supplies can be
moved between arbitrary leaves with positive probability. Replenishment in
the deleted leaf gives a self-loop. Thus the chain is irreducible and
aperiodic. Exchanging two sibling subtrees exchanges their counts and,
by \eqref{eq:policy}, their deletion masses. The transition kernel therefore
commutes with every tree automorphism. Uniqueness makes the stationary law
invariant, and the same kernel identity preserves any invariant initial law.
\end{proof}

\section{A multiscale drift estimate}\label{sec:energy}

The proof controls two state-dependent quantities: the terminal rate
energy \(H_L\) and the transport energy \(G\) from \eqref{eq:H-G}.
The potential below has a drift whose principal term dominates both.
Its curvature loss is at most \(LH_L\), which explains the choice
\(a\ge2000L\).

\subsection{The harmonic potential}

Define the harmonic inventory potential and its restoring term by
\begin{equation}\label{eq:Phi-D}
\Phi=\sum_{v\ne\mathrm{root}}p_v^2
\sum_{j=1}^{N_v}\frac1{j+a/2},
\qquad
D=\sum_{v\ne\mathrm{root}}p_v^2
\frac{p_v-q_v}{N_v+a/2}.
\end{equation}
Since \(t_vZ_v=(p_v-q_v)p_v/[a(N_v+a/2)]\),
\begin{equation}\label{eq:D-tZ}
D=a\sum_{v\ne\mathrm{root}}p_vt_vZ_v.
\end{equation}

\begin{lemma}[Exact harmonic drift]\label{lem:drift}
Conditional on the present count state,
\begin{equation}\label{eq:drift}
\E[\Phi_{\mathrm{next}}-\Phi\mid S]=D-\mathcal R,
\end{equation}
where
\begin{equation}\label{eq:curvature-error}
\mathcal R=
\sum_{v\ne\mathrm{root}}
\frac{p_v^3(1-q_v)}{(N_v+a/2)(N_v+a/2+1)}.
\end{equation}
Moreover,
\begin{equation}\label{eq:R-bound}
0\le\mathcal R\le\sum_{\ell=1}^L H_\ell\le LH_L.
\end{equation}
\end{lemma}
\begin{proof}
At a node of count \(k\), an unmatched replenishment increases its count
with probability \(p(1-q)\), while an unreplenished deletion decreases it
with probability \(q(1-p)\). Hence
\begin{align*}
\E\left[
\sum_{j=1}^{N_{\mathrm{next}}}\frac1{j+a/2}
-\sum_{j=1}^{k}\frac1{j+a/2}
\,\middle|\,S\right]
&=\frac{p(1-q)}{k+a/2+1}-\frac{q(1-p)}{k+a/2}\\
&=\frac{p-q}{k+a/2}
-\frac{p(1-q)}{(k+a/2)(k+a/2+1)}.
\end{align*}
This remains valid at \(k=0\). Multiplication by \(p_v^2\) and summation
give \eqref{eq:drift}.

Finally,
\[
\frac{p_v^3(1-q_v)}{(N_v+a/2)(N_v+a/2+1)}
\le p_vZ_v^2\le p_vh_v^2
\]
by \eqref{eq:invariant-domain}. Sum by levels and use
\eqref{eq:H-monotone}.
\end{proof}


\subsection{The local Bellman inequality}

Define the quadratic Bellman correction
\begin{equation}\label{eq:Bellman}
B(h,t,Z)=\left(t-\frac h2\right)Z-\frac{t^2}{6}.
\end{equation}
By Lemma~\ref{lem:invariants}, \(B\le0\) at every node.
For a quantity attached to the two children, write
\(\E_{\mathrm{ch}}\) for their equal average; these are spatial weights
\(1/2,1/2\), not deletion probabilities.

The mixed term \(-(h/2-t)Z\) couples the harmonic marginal \(Z\) to
the slack in the rate-floor constraint. The quadratic term uses the
identity \(t'_\pm=(t\mp b/a)/2\): its contribution to
\(\E_{\mathrm{ch}}B'-B\) is \(t^2/8-b^2/(24a^2)\).
Thus it supplies a positive squared-drift term at a controlled cost in
transport energy. The following inequality quantifies the resulting
balance.

\begin{lemma}[Local Bellman inequality]\label{lem:local}
At every internal node,
\begin{equation}\label{eq:local}
\E_{\mathrm{ch}}[t'Z'+B(h',t',Z')]-B(h,t,Z)
\ge\frac1{1000}\left\{
\E_{\mathrm{ch}}(h')^2-h^2+\frac{b^2}{a^2}
\right\}.
\end{equation}
\end{lemma}

The proof is a three-case calculation according to which term in
\eqref{eq:policy} is active. It uses only square completion,
Cauchy--Schwarz, Young's inequality, and elementary rational bounds. The
complete proof appears in Appendix~\ref{app:local-proof}.

\begin{proposition}[Deterministic aggregate estimate]\label{prop:aggregate}
Every count state satisfies
\begin{equation}\label{eq:aggregate}
D\ge
\frac a{1000}H_L+\frac{G}{1000a}
-\frac{501a}{1000m^2}.
\end{equation}
\end{proposition}
\begin{proof}
Multiply \eqref{eq:local} at node \(v\) by \(p_v\) and sum over all
internal nodes. Since each child has length \(p_v/2\), the left side
telescopes to
\[
\frac Da+\sum_{v\in\T_L}p_vB_v-B_{\mathrm{root}},
\]
while the right side telescopes to
\[
\frac1{1000}\left(H_L-H_0+\frac{G}{a^2}\right).
\]
The leaf terms satisfy \(B_v\le0\), and
\[
B_{\mathrm{root}}
=-\frac1{2m(m+a/2)}\ge-\frac1{2m^2}.
\]
Using \(H_0=m^{-2}\) gives \eqref{eq:aggregate}.
\end{proof}


\subsection{A unified energy bound}

\begin{proposition}[Energy estimate]\label{prop:energy}
Write \(H_L(t)\) and \(G(t)\) for the energies before period \(t\), and \(\Phi_t\) for the potential after that period. For every initial count law and every integer $T\ge1$,
\begin{equation}\label{eq:master-energy}
 \frac1T\sum_{t=1}^T
 \left\{\frac{a^2}{2}\E H_L(t)+\E G(t)\right\}
 \le \frac{501a^2}{m^2}
       +\frac{1000a}{T}\E[\Phi_T-\Phi_0].
\end{equation}
Consequently, under the stationary count law,
\begin{equation}\label{eq:stationary-energy}
 \E_\pi H_L\le\frac{1002}{m^2},\qquad
 \E_\pi G\le\frac{501a^2}{m^2}.
\end{equation}
For an arbitrary initial law,
\begin{equation}\label{eq:finite-G}
 \frac1T\sum_{t=1}^T\E G(t)
 \le\frac{501a^2}{m^2}+\frac{a^2}{2T}.
\end{equation}
\end{proposition}
\begin{proof}
Combining \eqref{eq:aggregate}, \eqref{eq:drift}, \eqref{eq:R-bound}, and $a\ge2000L$ gives
\[
 1000a\,\E[\Phi_{t}-\Phi_{t-1}\mid S_{t-1}]
 \ge\frac{a^2}{2}H_L(t)+G(t)-\frac{501a^2}{m^2}.
\]
Taking expectations and summing proves \eqref{eq:master-energy}.
Stationarity makes its final term zero. For the finite-time bound, use
\[
 0\le\Phi\le\log(1+2m/a)\le a/2000.
\]
The first upper bound follows by bounding each harmonic sum by its integral
and using $\sum_{v\ne\mathrm{root}}p_v^2<1$.
\end{proof}


\section{From energy to matching distance}\label{sec:transport}

For each internal node \(v\), let \(\tau_v\) be the tent supported on
\(I_v\), linear from zero at its left endpoint to height \(1/2\) at its
midpoint and back to zero at its right endpoint. If \(b_v\) uses actual
spatial labels, the deletion distribution has the exact integrated Haar
expansion
\begin{equation}\label{eq:Haar}
F_Q(z)-z=\sum_{v\in\T\setminus\T_L}b_v\tau_v(z),
\qquad
\int_0^1\tau_v(z)^2\,dz=\frac{p_v}{12}.
\end{equation}

\begin{lemma}[Cancellation and transport energy]\label{lem:transport-energy}
Under any tree-invariant count law,
\begin{equation}\label{eq:transport-energy}
\E\int_0^1(F_Q(z)-z)^2\,dz=\frac1{12}\E G.
\end{equation}
Consequently,
\begin{equation}\label{eq:transport-cost}
(\E C^2)^{1/2}\le\frac1n+\sqrt{\frac{\E G}{12}}.
\end{equation}
\end{lemma}
\begin{proof}
Tents belonging to disjoint nodes have disjoint interiors. For nested
distinct nodes, exchange the two child subtrees of the deeper node. Its
coefficient changes sign while the ancestor coefficient is fixed. Tree
invariance therefore makes their mixed expectation zero. Expanding
\eqref{eq:Haar} proves \eqref{eq:transport-energy}.


For the second assertion, write \(g(z)=F_Q(z)-z\). Since \(F_Q\) is
continuous and piecewise linear, \(F_Q(Y)=U\) almost surely, and
\[
\E[(Y-U)^2\mid S]
=\int_0^1g(z)^2F_Q'(z)\,dz
=\int_0^1g(z)^2\,dz+\left[\frac{g(z)^3}{3}\right]_0^1
=\int_0^1g(z)^2\,dz.
\]
The boundary term is zero. As \(|X-Y|\le1/n\), the \(L^2\) triangle
inequality and \eqref{eq:transport-energy} prove
\eqref{eq:transport-cost}.
\end{proof}


\subsection{Proof of Theorem~\ref{thm:main}}\label{sec:finite}

Under the stationary count law, \eqref{eq:stationary-energy} and
\eqref{eq:transport-cost} give
\begin{equation}\label{eq:stationary-cost}
(\E C_t^2)^{1/2}
\le\left(2+\sqrt{\frac{501}{12}}\right)\frac am
<8.5\frac am.
\end{equation}
For an arbitrary initial inventory, the finite count chain converges to
its stationary law by Lemma~\ref{lem:symmetry}. In particular, the
expectation of the bounded count function
\(\int_0^1(F_Q(z)-z)^2\,dz\) converges to its stationary expectation.
The conditional second-moment identity in the proof of
Lemma~\ref{lem:transport-energy} and the bound \(|X-Y|\le1/n\)
therefore give the same limiting bound as \eqref{eq:stationary-cost}.
Since \(a=O(\log m)\), this proves part (i).

If the initial count law is tree invariant, it remains so at every time.
Apply \eqref{eq:transport-cost} at a uniformly selected time among
\(1,\ldots,S\), and use \eqref{eq:finite-G}, to obtain
\[
\left(\frac1S\sum_{t=1}^S\E C_t^2\right)^{1/2}
\le \frac{2a}{m}
+\sqrt{\frac1{12}\left(\frac{501a^2}{m^2}+\frac{a^2}{2S}\right)}.
\]
For \(S\ge m^2\), this is less than \(8.5a/m\) and also bounds
the average expected distance.
From an arbitrary initial inventory, replace one original supply per
period during the first \(m\) periods. This costs at most \(m\) and
leaves iid uniform supplies. For \(T\ge2m^2\), the remaining horizon
\(S=T-m\) is at least \(m^2\), so
\[
J_T\le\frac mT+\frac ST\,8.5\frac am<9\frac am.
\]
This proves part (ii). The implementation bounds were established in
Section~\ref{sec:policy}.
\hfill\(\square\)

\begin{corollary}[Balanced initial inventory]\label{cor:balanced}
For every \(m\), there is an initial inventory law for which
\[
\left(\frac1T\sum_{t=1}^T\E C_t^2\right)^{1/2}
\le\left(2+\sqrt{\frac{501}{12}}\right)\frac am
\qquad\text{for every }T\ge1.
\]
When \(n\) divides \(m\), equal counts in all leaves suffice.
\end{corollary}
\begin{proof}
Split the \(m\) supplies recursively as evenly as possible, using
independent fair coins to choose the larger child whenever a count is odd.
At each level the counts differ by at most one. Concavity of the harmonic
summand in \eqref{eq:Phi-D} implies that this configuration maximizes
the contribution of every level, hence maximizes \(\Phi\).
Its count law is tree invariant. Thus \(\Phi_T-\Phi_0\le0\), and
\eqref{eq:master-energy} bounds the average of \(\E G(t)\) by
\(501a^2/m^2\) for every \(T\). Apply
\eqref{eq:transport-cost} at a uniformly selected time.
Positions within the leaves are arbitrary.
\end{proof}

\section{Design principles}\label{sec:heuristic}

The preceding proof is exact. We now give complementary, near-balance
calculations that explain the shape of the policy and potential.
The approximations in this section are not used in the proof.

\subsection{The logarithm and its spatial weight}

An empty interval of length \(p\) receives demands with probability \(p\),
and serving them costs a distance of order \(p\). Even if supplies are
placed just outside both endpoints, its cost contribution is
\(\int_0^p\min\{z,p-z\}\,dz=p^2/4\).
This suggests the weight \(p^2\) in \(\Phi\).
The harmonic summand is a discrete, regularized logarithm: its exact
marginal deletion loss is \(1/(N+a/2)\), approximately \(1/N\) when
\(N\gg a\).

The logarithm measures scarcity relative to the typical inventory \(mp\).
Writing \(r=N/(mp)\), we have
\[
\log N=\log(mp)+\log r,\qquad
\frac{d^2}{dr^2}\log r=-\frac1{r^2}\approx-1\quad(r\approx1).
\]
Thus it has order-one curvature in relative-inventory coordinates at
every spatial scale. Near balance, a unit inventory change produces a
relative change of order \(1/(mp)\) and occurs with probability of order
\(p\). Its contribution to the curvature loss is consequently
\[
p^2\cdot p\cdot\frac1{(mp)^2}=\frac p{m^2}
\]
per node, or \(1/m^2\) per level. The exact counterpart is
\(\mathcal R\le LH_L\).

\subsection{The feedback as penalized potential ascent}

For positive sibling counts \(x\ge y\), the first-order deletion
contribution to \(\log x+\log y\) is
\[
-\frac{q+b}{2x}-\frac{q-b}{2y}
=-\frac q2\left(\frac1x+\frac1y\right)
+\frac b2\left(\frac1y-\frac1x\right).
\]
The term depending on \(b\) favors deletion from the larger inventory.
Let \(\bar b\) be the smaller of the two caps in \eqref{eq:policy}.
The policy is exactly the maximizer of
\begin{equation}\label{eq:penalized-ascent}
\max_{0\le b\le\bar b}
\left\{\frac b2\left(\frac1y-\frac1x\right)-\frac{b^2}{8aq}\right\}.
\end{equation}
The unconstrained maximizer is \(2aq(x-y)/(xy)\). This is a local
logarithmic approximation to the drift design; it is not an exact
optimization of the full, discrete potential \(\Phi\).

The penalty must be compared with the spatial weight of the potential.
Both children have weight \(p^2/4\), so the weighted penalty is
\[
\frac{p^2b^2}{32aq}
=\frac{3p}{8aq}\left(\frac{pb^2}{12}\right).
\]
When \(q\approx p\), this charges the same multiple of squared transport
energy at every scale. Equivalently, with
\(\rho=(x-y)/(x+y)\) and \(\beta=b/q\), the unweighted objective is
approximately
\[
\frac2m\frac{\rho\beta}{1-\rho^2}-\frac p{8a}\beta^2.
\]
The potential benefit is scale free in these coordinates, while the
penalty for a fixed relative bias increases with interval length.

\subsection{Why comparable fluctuations across levels?}

For any stationary policy, using fixed left--right labels and
\(\Delta=N_{\mathrm L}-N_{\mathrm R}\), one has the exact identities
\begin{equation}\label{eq:stationary-heuristic}
\E q=p,\qquad \E[\Delta b]=p.
\end{equation}
The first is inventory conservation. For the second, let \(A,B\) be the
signed child indicators of replenishment and deletion. Then
\(\Delta'=\Delta+A-B\), with conditional moments
\(\E A=0\), \(\E A^2=p\), \(\E B=b\), \(\E B^2=q\), and
\(\E AB=0\). The drift of \(\Delta^2\) is therefore
\(-2\Delta b+p+q\). Stationarity gives \eqref{eq:stationary-heuristic}.

At a well-stocked node, approximate \(N\) by \(mp\) and the feedback by
\(b=k_p\Delta\) with deterministic \(k_p\). Let
\(s_\ell=\E(\Delta/N)^2\) for a node at level \(\ell\) under the
tree-invariant stationary law, with \(\Delta/N=0\) when \(N=0\). Equation
\eqref{eq:stationary-heuristic} then predicts
\[
s_\ell\approx\frac1{k_pm^2p},\qquad
E_\ell:=\frac1{12}\sum_{v\in\T_\ell}p_v\E b_v^2
\approx\frac1{12m^2s_\ell}.
\]
Allowing more relative fluctuation makes a level cheaper to stabilize.
The logarithm suggests how these fluctuations accumulate. When every
leaf is occupied, set \(r_v=N_v/(mp_v)\) and
\(\rho_v=(N_{v_{\mathrm L}}-N_{v_{\mathrm R}})/N_v\).
The identities \(r_{v_{\mathrm L}}=r_v(1+\rho_v)\) and
\(r_{v_{\mathrm R}}=r_v(1-\rho_v)\) telescope to
\[
-\sum_{i\in\T_L}p_i\log r_i
=\frac12\sum_{v\in\T\setminus\T_L}p_v[-\log(1-\rho_v^2)]
.
\]
Expanding \(-\log(1-\rho^2)\) near zero suggests an expected logarithmic
scarcity of approximately \(\frac12\sum_\ell s_\ell\). Under a
provisional bound \(\sum_\ell s_\ell\le C_0\) on accumulated scarcity,
Cauchy--Schwarz gives
\(\sum_\ell1/s_\ell\ge L^2/C_0\), with equality at equal
\(s_\ell\). This motivates distributing relative-inventory risk
comparably across well-stocked levels.

It also recovers the potential weight. A local penalty
\(-w(p)\rho^2\approx-w(p)\Delta^2/(m^2p^2)\) has derivative
\(-2w(p)\Delta/(m^2p^2)\) with respect to \(\Delta\). Since bias \(b\)
reduces the expected imbalance by \(b\), its first-order benefit is
\(2w(p)b\Delta/(m^2p^2)\). Charging \(\eta pb^2\) for transport
gives \(b=w(p)\Delta/(\eta m^2p^3)\). Comparable risk across levels
requires a feedback coefficient proportional to \(1/p\), hence
\(w(p)\propto p^2\).
For \eqref{eq:policy}, the interior prediction is
\(b\approx8a\Delta/(m^2p)\), giving \(s_\ell\approx1/(8a)\)
and \(E_\ell\approx2a/(3m^2)\). These scale profiles remain heuristic;
the proof bounds their aggregate and handles depletion without a
near-balance assumption.

\section{Further directions}\label{sec:extensions}

The drift calculation uses the nested count structure, conservation, and
bounded inventory increments. These ingredients can arise in other
inventory and stochastic allocation systems. The transport step is more
specific: monotone coupling on the line and symmetry of the binary
partition turn deletion imbalances into an additive squared-distance
bound. Extending the method to another geometry requires a corresponding
transport-energy estimate; the local Lyapunov calculation alone does not
provide one.

\paragraph{Connection to Shor's Problem D.}
A spatial sweep suggests a link with Shor's rightward matching problem
in the unit square \cite{Shor}. Sweep from right to left, use vertical
coordinates as one-dimensional locations, and hold encountered red
points in a replenishment queue. Fictitious supplies cover queue
shortages and correspond to matches to the bottom boundary.
Combined with a balanced-inventory version of our bound, this gives a
route to expected total length of order
\[
\frac{N\log(m+1)}m+m+\sqrt N.
\]
Choosing \(m\asymp\sqrt{N\log(N+1)}\) would recover the order of Shor's
upper bound. We leave the reduction and its queue accounting outside the
scope of this paper.

\appendix
\section{Proof of the local Bellman inequality}\label{app:local-proof}

We prove Lemma~\ref{lem:local}. First suppose both child counts are
positive and temporarily order them as \(x\ge y\). Normalize all rates by
the parent rate \(h=q/N\), and set
\begin{equation}\label{eq:normalized}
s=\frac Na,\quad r=\frac th,\quad \rho=\frac{x-y}{N},\quad
\lambda=\rho^2,\quad A=s(1-\lambda),\quad
g=s+r=\frac p{ah},\quad z=\frac b{ah}.
\end{equation}
The domain is \(s>0\), \(-s<r\le1/2\), and \(0\le\rho<1\).
Writing hats for division by \(h\), direct calculation gives
\begin{equation}\label{eq:child-normalized}
\widehat h_\pm=\frac{s\pm z}{s(1\pm\rho)},\qquad
\widehat t_\pm=\frac{r\mp z}{2},\qquad
\widehat Z_\pm=\frac{g}{s(1\pm\rho)+1},\qquad
\frac Zh=\frac{g}{s+1/2}.
\end{equation}
The three candidates in \eqref{eq:policy}, divided by \(ah\), become
\begin{equation}\label{eq:normalized-policy}
z_F=\frac{8s\rho}{A},\qquad
z_U=s\rho,\qquad
z_H=\frac{s-s(1-\rho)r}{s(1-\rho)+1}.
\end{equation}

Assume first that \(\rho>0\), put
\begin{equation}\label{eq:kJd}
k=\frac z\rho,\qquad J=\frac{s-k}{A},\qquad d=s(A+2)+1,
\end{equation}
and let \(R_0\) denote the left side of \eqref{eq:local}, divided by
\(h^2\). Expanding \eqref{eq:Bellman} gives the common identity
\begin{align}
R_0={}&\frac{r^2}{8}-\frac{\lambda k^2}{24}
+\frac{g(1/2-r)(s+1)}{2d(s+1/2)} \notag\\
&+\frac{\lambda g}{d}
\left[sk-J(s+1/2)-\frac{(1/2-r)s^2}{s+1/2}\right].
\label{eq:R0}
\end{align}
One way to verify \eqref{eq:R0} is first to calculate
\[
\frac{\widehat Z_++\widehat Z_-}{2}=\frac{g(s+1)}d,
\quad
\frac{\widehat h_++\widehat h_-}{2}=1+\lambda J,
\quad
\frac{\widehat h_- -\widehat h_+}{2}=\rho J,
\]
and then expand the three terms of \(B\). The child rate-energy increment is
\begin{equation}\label{eq:rate-increment}
\E_{\mathrm{ch}}\widehat h'^2-1
=2\lambda J+(1+\lambda)\lambda J^2.
\end{equation}
If \(\rho=0\), then \(z=0\), the right side of \eqref{eq:local} is zero,
and direct substitution leaves the nonnegative first and third terms of
\eqref{eq:R0}. We henceforth take \(\rho>0\). Ties between candidates may
be assigned to either applicable case.

\subsection{The feedback is active}

Here \(k=8s/A\). Comparison with the uniform-deletion cap gives
\(A\ge8\), and hence \(s\ge A\ge8\), while
\[
J=\frac{s(A-8)}{A^2}\ge0.
\]
Using \(r\le1/2\) and \(s^2/(s+1/2)\le s\), the bracket in
\eqref{eq:R0} is at least
\begin{equation}\label{eq:feedback-bracket}
rs+\frac{7s^2}{A}+\frac{8s^2}{A^2}
-\frac s2-\frac{s}{2A}
\ge rs+\frac{13s^2}{2A}.
\end{equation}
For the second inequality, use \(s/2\le s^2/(2A)\) and
\(8(s/A)^2-s/(2A)\ge0\). The first bound has also discarded a positive
term \(4s/A^2\).

Set
\[
\omega=\frac sA\ge1,\qquad
u=\frac r\omega,\qquad
\alpha=\frac1A\le\frac18,\qquad
C=\frac{13}{2},\qquad
\delta=\frac d{sA}.
\]
Then
\[
1+2\alpha\le\delta\le(1+\alpha)^2,
\qquad
g=s(1+\alpha u)>0.
\]
Discarding the nonnegative third term in \eqref{eq:R0} and applying
\eqref{eq:feedback-bracket} gives
\[
\frac{R_0}{\omega^2}
\ge\frac{u^2}{8}
+\frac{\lambda}{\delta}(1+\alpha u)(u+C)
-\frac83\lambda.
\]
Completing the square in \(u\) bounds this below by
\begin{equation}\label{eq:feedback-square}
\lambda\left[
\frac C\delta-\frac83-
\frac{2\lambda(1+C\alpha)^2}
{\delta(\delta+8\lambda\alpha)}
\right].
\end{equation}
Since \(0\le\lambda\le1\),
\[
\frac{2\lambda(1+C\alpha)^2}
{\delta(\delta+8\lambda\alpha)}
\le\frac{2(1+13\alpha/2)^2}{(1+2\alpha)(1+10\alpha)}
\le\frac{19}{8}.
\]
For the last inequality, cross multiplication leaves
\[
\frac38+\frac52\alpha-37\alpha^2
\ge\frac38-\frac{17}{8}\alpha\ge\frac7{64},
\]
because \(0\le\alpha\le1/8\). Also
\(C/\delta\ge416/81>41/8\). Therefore
\begin{equation}\label{eq:feedback-R}
R_0\ge\left(\frac{41}{8}-\frac83-\frac{19}{8}\right)
\lambda\omega^2=\frac1{12}\lambda\omega^2.
\end{equation}
On the other hand, \(0\le J\le\omega\) and
\(z^2=64\lambda\omega^2\). Equation~\eqref{eq:rate-increment} gives
\[
\E_{\mathrm{ch}}\widehat h'^2-1+z^2
\le68\lambda\omega^2.
\]
Since \(1/12>68/1000\), this proves \eqref{eq:local} in the feedback case.

\subsection{The uniform-deletion cap is active}

Here \(z=s\rho\), so both child rates equal \(h\) and the rate-energy
increment is zero. Comparison with the feedback gives \(A\le8\).
Equation~\eqref{eq:R0} simplifies to
\begin{equation}\label{eq:uniform-R}
R_0=\frac{r^2}{8}-\frac{z^2}{24}
+\frac{g^2z^2}{d(s+1/2)}
+\frac{g(1/2-r)(s+1)}{2d(s+1/2)}.
\end{equation}

If \(s\ge1/2\), the bounds \(A\le\min\{s,8\}\) give
\[
\frac{d(s+1/2)}{s^2}
=A+2+\frac{A}{2s}+\frac2s+\frac1{2s^2}
\le
\begin{cases}
19/2,&1/2\le s\le1,\\
13,&s\ge1.
\end{cases}
\]
Thus \(d(s+1/2)\le13s^2\). Since \(g-r=s\), weighted
Cauchy--Schwarz gives
\[
\frac{r^2}{8}+\frac{\lambda g^2}{13}
\ge\frac{\lambda(g-r)^2}{13+8\lambda}
\ge\frac{z^2}{21}.
\]
The last term in \eqref{eq:uniform-R} is nonnegative, and hence
\[
R_0\ge z^2\left(\frac1{21}-\frac1{24}\right)
=\frac{z^2}{168}>\frac{z^2}{1000}.
\]

If \(s\le1/2\), use \(d\le(s+1)^2\) to bound the last term in
\eqref{eq:uniform-R} below by \(g(1/2-r)/3\). Discarding the positive
squared term and using \(z^2\le s^2\) gives
\[
R_0\ge
\frac{r^2}{8}-\frac{s^2}{24}
+\frac{(s+r)(1/2-r)}3.
\]
The right side is concave in \(r\in[-s,1/2]\). Its endpoint values are
\(s^2/12\) and \(1/32-s^2/24\ge s^2/12\). Thus
\(R_0\ge z^2/12\), completing this case.

\subsection{The rate-floor cap is active}

Use coordinates adapted to the active equality:
\begin{equation}\label{eq:floor-coordinates}
u=\frac{2y}{a},\qquad
K=\frac{2h_-}{h}-1,\qquad
e=1-\frac{2t}{h},\qquad
v=\frac{2x}{a}.
\end{equation}
The uniform-deletion cap gives \(K\ge1\), and
\eqref{eq:invariant-domain} gives \(e\ge0\). The rate-floor equality
implies
\begin{equation}\label{eq:floor-identities}
v=K(u+1)+e,\qquad
z=\frac{K+e}{2},\qquad
\widehat h_- =\widehat Z_- =\frac{K+1}{2}.
\end{equation}
Keep the larger child's normalized rate and inverse inventory in
\[
\ell=\widehat h_+
=1-\frac{u(K-1)}{2v},\qquad
w=\widehat Z_+
=\frac{(K+1)(u+1)}{2(v+1)},
\]
and put
\[
\eta=\frac{v+1-u}{v+1+u}.
\]
The normalized drifts are
\[
\widehat t_+=-\frac{2e+K-1}{4},\qquad
\widehat t_- =\frac{K+1}{4},
\]
while \(Z/h=(K+1)(u+1)/(u+v+1)\). Substitution in
\eqref{eq:Bellman} gives
\begin{equation}\label{eq:floor-R}
R_0=
\frac{5K^2+12K+9+2e^2-2Ke-6e}{96}
+\frac w2\left(e\eta-\frac{K-1+\ell}{2}\right).
\end{equation}

We use only the following three consequences of the policy:
\begin{equation}\label{eq:floor-bounds}
0<\ell\le1,\qquad
0<w\le\frac{K+1}{2K},\qquad
\eta\ge\frac{K+e}{16+K+e}.
\end{equation}
The first two follow from \(v=K(u+1)+e\). For the third, comparison with
the feedback candidate gives
\[
4(v^2-u^2)\ge(K+e)uv.
\]
Put \(\theta=u/v\) and \(c=(1-\theta)/(1+\theta)\). Then
\[
16c\ge(K+e)(1-c^2)\ge(K+e)(1-c),
\]
so \(c\ge(K+e)/(16+K+e)\); since \(\eta\ge c\), the final bound in
\eqref{eq:floor-bounds} follows.

If \(e\eta\ge K/2\), the second term in \eqref{eq:floor-R} is
nonnegative. Using \(2(K+3)e\le e^2+(K+3)^2\) in the first term gives
\begin{equation}\label{eq:floor-first-case}
R_0\ge\frac{4K^2+6K+e^2}{96}.
\end{equation}
If \(e\eta<K/2\), use \(\ell\le1\), then the upper bound on \(w\), and
finally the lower bound on \(\eta\). Because the intervening factor is
negative, the inequality directions are preserved. The result is
\begin{equation}\label{eq:J-bound}
96R_0\ge J(K,e),
\end{equation}
where
\[
J(K,e)=5K^2-3+2e^2-(2K+6)e
+\frac{24(K+1)e(K+e)}{K(16+K+e)}.
\]
We claim that
\begin{equation}\label{eq:J-claim}
J(K,e)\ge\frac{K^2+e^2}{4}
\qquad(K\ge1,\ e\ge0).
\end{equation}
To prove \eqref{eq:J-claim}, first suppose \(e\ge K\). Then
\[
\frac{24(K+1)(K+e)}{K(16+K+e)}
\ge\frac{24(K+1)}{K+8},
\qquad
2K+6-\frac{24(K+1)}{K+8}\le\frac{8K}{3}.
\]
The second inequality follows by cross multiplication from
\(2(K-1)(K+36)\ge0\). Hence
\[
J(K,e)\ge2K^2+2e^2-\frac83Ke
\ge\frac23(K^2+e^2).
\]
Here we used \(K\ge1\) and \(2Ke\le K^2+e^2\).

Now suppose \(e\le K\), and put \(\xi=e/K\in[0,1]\) and
\(c=1+\xi\in[1,2]\). Since
\[
\frac{K+1}{16+Kc}\ge\frac2{16+c},
\qquad
6-\frac{48c}{16+c}\ge\frac23,
\]
we have
\begin{align*}
\frac{J(K,e)}{K^2}
&\ge5+2\xi^2-2\xi-\frac3{K^2}
 -\frac{\xi}{K}\left(6-\frac{48c}{16+c}\right)\\
&\ge2+2\xi^2-8\xi+\frac{48\xi(1+\xi)}{17+\xi}
=:j(\xi).
\end{align*}
The bound \((K+1)/(16+Kc)\ge2/(16+c)\) is equivalent to
\((K-1)(16-c)\ge0\); the second line uses \(K\ge1\).
Finally,
\[
4(17+\xi)\left[j(\xi)-\frac{1+\xi^2}{4}\right]
=7\xi^3+29(3\xi-2)^2+18\xi^2+3\xi+3\ge0.
\]
This proves \eqref{eq:J-claim} in both cases.


Equations~\eqref{eq:floor-first-case} and \eqref{eq:J-claim} imply in
the two cases that
\[
R_0\ge\frac{K^2+e^2}{384}.
\]
Moreover, \(\ell\le1\), \((K+1)^2/4\le K^2\), and
\((K+e)^2/4\le(K^2+e^2)/2\) give
\[
\E_{\mathrm{ch}}\widehat h'^2-1+z^2
\le\frac{(K+1)^2/4-1}{2}+\frac{(K+e)^2}{4}
\le K^2+e^2.
\]
Since \(1/384>1/1000\), this proves \eqref{eq:local} in the rate-floor
case.

\subsection{Empty-node boundaries}

When the smaller child becomes empty, hold the parent variables fixed and
let its count tend continuously to zero in the preceding formulas. The
feedback candidate diverges. If \(g=p/(ah)<1\), the uniform-deletion cap
is eventually active and the limiting auxiliary empty-child rate is
\(h\). If \(g>1\), the rate-floor cap is eventually active and the limiting
rate is \(gh=p/a\). At \(g=1\) the limits agree. This is precisely the
definition \eqref{eq:empty-rate}, so \eqref{eq:local} extends continuously.
Equivalently, the floor formulas apply directly at \(u=0\), where
\(v=K+e\), \(g=(K+1)/2\ge1\), and the occupied child's normalized rate
equals one.

At an empty parent, \(t=p/a\), \(Z=2t\), and both children inherit \(h\),
with \(t'=t/2\) and \(Z'=t\). Direct calculation gives
\[
\E_{\mathrm{ch}}[t'Z'+B']-B
=\frac{ht}{2}-\frac{7t^2}{8}
\ge\frac{ht}{16},
\]
because \(0<t\le h/2\). Both terms on the right side of
\eqref{eq:local} vanish. This completes the proof of
Lemma~\ref{lem:local}.


\section*{AI assistance}
OpenAI's ChatGPT and Codex were used in developing the policy and proof,
checking algebra, and drafting and revising the manuscript.

\section*{Data availability}
No datasets were generated or analyzed. All proofs are contained in the
article.

\begin{thebibliography}{9}

\bibitem{AKT}
M.~Ajtai, J.~Koml\'os, and G.~Tusn\'ady.
On optimal matchings.
\emph{Combinatorica}, 4(4):259--264, 1984.
\href{https://doi.org/10.1007/BF02579135}{doi:10.1007/BF02579135}.

\bibitem{BFP}
E.~Balkanski, Y.~Faenza, and N.~P\'erivier.
The power of greedy for online minimum cost matching on the line.
In \emph{Proceedings of the 24th ACM Conference on Economics and
Computation}, pages 185--205, 2023.
Full version: \href{https://arxiv.org/abs/2210.03166v3}{arXiv:2210.03166v3},
2025.

\bibitem{HPZ}
N.~Holden, Y.~Peres, and A.~Zhai.
Gravitational allocation for uniform points on the sphere.
\emph{The Annals of Probability}, 49(1):287--321, 2021.
\href{https://doi.org/10.1214/20-AOP1452}{doi:10.1214/20-AOP1452}.

\bibitem{Kanoria}
Y.~Kanoria.
Dynamic spatial matching.
\emph{arXiv:2105.07329v3}, 2023.
\href{https://arxiv.org/abs/2105.07329v3}{arXiv:2105.07329}.

\bibitem{Shor}
P.~W.~Shor.
How to pack better than Best Fit: tight bounds for average-case online
bin packing.
In \emph{Proceedings of the 32nd Annual Symposium on Foundations of
Computer Science}, pages 752--759, 1991.
\href{https://doi.org/10.1109/SFCS.1991.185444}{doi:10.1109/SFCS.1991.185444}.

\end{thebibliography}
\end{document}
